% C44_POS /root/axis_independent; GPT-6 Astra Pro; effort setting not exposed.
% Mathematical deductions for root review; no claim of formal verification.
\section{The exact axis primitive space and the limits of the proposed transfers}

\paragraph{Notation and carried input.}
Put $r=m-1$, so that $\dim X=2r+2$ and $\dim B=2r$.
Write $L_X$ and $L_B$ for multiplication by $h$ and $h_B$.
All cycle groups have rational coefficients, and $\operatorname{Alg}^j$
means the image of the rational cycle-class map for one fixed
$\ell\ne p$. No equality between numerical and homological equivalence
on $B$ is assumed. The carried statement JAC1964 H9-T7, equation
(H9.18), is the isometry
\[
 A_X\oplus A_B\ \simeq\ \mathcal H/\mathcal H_0,
 \qquad
 Q((z,w),(z',w'))=(-1)^m\int_Xzz'+(-1)^r\int_Bww'.
\]
Its proof establishes $L_Bw=0$ for $w\in A_B$, not merely
$L_B^2w=0$. The factorization $B\hookrightarrow Y\hookrightarrow X$
through a smooth pencil member and the injectivity of the latter
Gysin map are essential to that stronger conclusion. We use this
checked input without reproving its height isometry.

\subsection*{C44-A01. The axis gate is the entire remaining primitive sign on the axis}

\textbf{Proposition.}
There is an orthogonal direct sum of rational algebraic classes
\[
 \operatorname{Alg}^{r}(B)\cap P^{2r}(B)
 =\kappa^*\bigl(\operatorname{Alg}^{r}(X)\cap P^{2r}(X)\bigr)
       \mathbin{\perp} A_B,                                  \tag{C44.A01.1}
\]
where $P^{2r}(B)=\ker L_B$ and $P^{2r}(X)=\ker L_X^3$.
On the first summand the signed middle form is
\[
   (-1)^r\int_B\kappa^*a\,\kappa^*a'
       =(-1)^r\int_Xh^2aa'.                                 \tag{C44.A01.2}
\]
For every $0\le j<r$, restriction identifies the primitive algebraic
classes of codimension $j$ on $X$ and $B$, including their signed
Lefschetz pairings.
Consequently, \emph{conditional on the unrefereed 032/Milne input
asserting the Hodge standard conjecture for $X$}, strict positivity
on $A_B$ is equivalent to the full Hodge standard conjecture $I(B)$
for $h_B$. Weak nonnegativity on $A_B$ is equivalent to the version
of $I(B)$ with every positive-definite assertion replaced by
positive-semidefinite. These last two statements must not be
identified unless nondegeneracy on algebraic classes is additionally
known.

\textbf{Proof.}
Hard Lefschetz makes $L_X^2:H^{2r}(X)\longrightarrow H^{2r+4}(X)$
an isomorphism. Its inverse $\beta$ is induced by an algebraic
correspondence because the Lefschetz standard conjecture is known
for abelian varieties, without 032. The self-intersection and
projection formulas give
$\kappa_*\kappa^*=L_X^2$. Hence
$e=\kappa^*\beta\kappa_*$ is a projector and every algebraic middle
class $x$ of $B$ has a unique decomposition
\[
 x=\kappa^*a+w,\qquad a=\beta\kappa_*x,\qquad w\in A_B.       \tag{C44.A01.3}
\]
Both components are rational algebraic classes. The summands are
orthogonal because
$\int_B\kappa^*a\,w=\int_Xa\,\kappa_*w=0$.
H9-T7 gives $L_Bw=0$. The Gysin map
\[
 \kappa_*:H^{2r+2}(B)\longrightarrow H^{2r+6}(X)
\]
is an isomorphism, being dual to weak Lefschetz in degree $2r-2$.
Therefore $L_B\kappa^*a=0$ if and only if
$\kappa_*L_B\kappa^*a=L_X^3a=0$. This proves (C44.A01.1).
The projection formula proves (C44.A01.2).

For $j<r$, weak Lefschetz gives an isomorphism
$\kappa^*:H^{2j}(X)\longrightarrow H^{2j}(B)$.
Its inverse on this degree is the algebraic operator
\[
 (L_X^{2r+2-2j})^{-1}\kappa_*L_B^{2r-2j};                   \tag{C44.A01.4}
\]
indeed its composition with $\kappa^*$ is the identity.
Thus this is an isomorphism on rational algebraic classes too.
For $1\le j<r$, the map $\kappa_*$ in the degree of
$L_B^{2r-2j+1}\kappa^*a$ is an isomorphism by dual weak
Lefschetz in degree $2j-2$. Thus the primitive condition on $B$
is equivalent to
\[
 L_X^{2r-2j+3}a=0,
\]
which is precisely the primitive condition on $X$. For $j=0$
both primitive spaces are $H^0$, so the same assertion holds.
Finally
\[
 (-1)^j\int_B h_B^{2r-2j}\kappa^*a\,\kappa^*a'
 =(-1)^j\int_Xh^{2r+2-2j}aa'.
\]
The conditional positivity assertions now follow from the
orthogonal direct sum and the signed Hodge standard conjecture on
$X$. No positivity assertion on the new summand has been proved.
\hfill$\square$

\paragraph{Programme location and scope.}
This is a sharpened identification of JAC1964 QH9-axis, and concerns
cote 44, Th.~5, scan 004, with the primitive margin on scan 008.
In cote 16 terminology it isolates $I(B)$; the algebraic operator
$\beta$ uses $C_G(X)=B(X)$, not $I(B)$. It does not close
QH9-axis and does not upgrade its requested semidefinite sign to
strict positivity.

\subsection*{C44-A02. Every graph transfer to an abelian target annihilates the axis classes}

\textbf{Proposition.}
Let $D$ be an abelian variety and $f:B\longrightarrow D$ any
morphism. Then $f_*w=0$ in cohomology for every $w\in A_B$.
The same holds for a rational map, interpreted through its unique
extension. Consequently $A_B$ is orthogonal to every cohomology
class pulled back by such a morphism in the complementary degree.
This is unconditional. It does not rule out correspondences on
$B\times D$ which are not assembled from these graph maps.

\textbf{Proof.}
Choose a geometric base point $b$ on $B$. Let
$a_B:B\longrightarrow\operatorname{Alb}(B)$ be the pointed Albanese
map and let $u:\operatorname{Alb}(B)\longrightarrow X$ be induced by
$x\mapsto\kappa(x)-\kappa(b)$.
Weak Lefschetz in degree one identifies $H^1(X)$ with $H^1(B)$.
The identification of degree-one cohomology with the dual of the
Albanese Tate module therefore makes $u$ an isogeny.
After subtracting the constant $f(b)$, the morphism $f$ factors as
$v a_B$. Since $u$ is invertible in the isogeny category, there are
an integer $N>0$ and a homomorphism $t:X\longrightarrow D$ such that
$t u=[N]_D v$. Thus
\[
 [N]_D\circ f=t\circ\kappa+\text{a constant}.
\]
Translations on an abelian variety act trivially on cohomology.
Functoriality of proper pushforward gives
\[
 [N]_{D,*}f_*w=t_*\kappa_*w=0.
\]
Multiplication by $N$ induces an invertible map on rational
$\ell$-adic cohomology, so $f_*w=0$. The orthogonality assertion is
projection formula. A rational map from a smooth variety to an
abelian variety extends by Weil's extension theorem; one exact
modern locator is Milne, \emph{Abelian Varieties}, course notes
version 2.00, I, Theorem 3.2, printed p.~17 (PDF p.~23).\hfill$\square$

\paragraph{Programme location and scope.}
This specifies a limit of the proposed curve-test attack on
QH9-axis. It concerns the incidence higher Picard framework of
cote 44, Th.~1--4, scan 002, and the incidence divisor on scans
008--009. It complements JAC-P19 (no abelian rational domination)
and the already proved annihilation calculation in H9-T13, rather
than replacing either. An algebraic correspondence shifts
cohomological degree by an even number, so $H^1(C)$ by itself
cannot map into $H^{2r}(B)$ through an algebraic correspondence.
The elliptic tensor construction $w\boxtimes P_E$ resolves that
parity issue but its ambient pushforward is zero, exactly as
recorded in H9-T13. The missing input is a genuinely new
pairing-compatible correspondence, not a graph of another
Albanese map.

\subsection*{C44-A03. The finite-field test is absence of potential Tate roots, not one exceptional slope}

\textbf{Proposition.}
Suppose the axis and the embedding are defined over $\mathbf F_q$.
Let
\[
 V_B=\ker\bigl(\kappa_*:H^{2r}(B_{\overline{\mathbf F}_q},
 \mathbf Q_\ell)\longrightarrow H^{2r+4}(X_{\overline{\mathbf F}_q},
 \mathbf Q_\ell)(2)\bigr)
\]
be the untwisted primitive variable cohomology, with the twists in
the target understood by Gysin compatibility. Write $F$ for
geometric Frobenius and $P_B(t)=\det(t-F\mid V_B)$.
Every algebraic axis class is contained in
\[
 T_B=\bigcup_{a\ge1}\ker(F^a-q^{ra}\mid V_B).                \tag{C44.A03.1}
\]
In particular all its normalized Newton slopes are $r$.
If $P_B$ has no root $q^r\zeta$ for a root of unity $\zeta$,
then $A_B=0$.
For $D=\deg P_B>0$, a finite exact sufficient test is
\[
 \gcd\bigl(P_B(t),\ t^N-q^{rN}\bigr)=1,
 \qquad N=\operatorname{lcm}(1,2,\ldots,2D^2).               \tag{C44.A03.2}
\]
Here $P_B\in\mathbf Q[t]$ is obtained by dividing the middle
Frobenius polynomial of $B$ by that of the restricted ambient
summand. No Tate conjecture and no Frobenius semisimplicity
conjecture is needed for this vanishing criterion.

\textbf{Proof.}
A rational algebraic cycle on $B_{\overline{\mathbf F}_q}$ is
defined over a finite extension $\mathbf F_{q^a}$. Its codimension
$r$ cycle class in untwisted cohomology satisfies $F^a w=q^{ra}w$.
This proves (C44.A03.1). In crystalline realization, the same
cycle is a Frobenius eigenclass of slope $r$; equivalently every
root $q^r\zeta$ has normalized valuation $r$ at every prime over
$p$.
If the polynomial has no such root, every kernel in (C44.A03.1)
is zero, which proves the asserted vanishing.

The splitting (C44.A01.3) is Frobenius-compatible after the
polarization is defined over $\mathbf F_q$; its image is isomorphic
to $H^{2r}(X)$. Thus the usual rational Frobenius polynomials of
$X$ and $B$ give $P_B\in\mathbf Q[t]$ without asserting
rationality of a chosen algebraic basis of $A_B$.
If $q^r\zeta$ is a root and $n$ is the order of $\zeta$, the
cyclotomic polynomial of $\zeta$ divides
$q^{-rD}P_B(q^rt)$, so $\varphi(n)\le D$.
For $n=\prod p^a$,
\[
 \frac{\varphi(n)^2}{n}=\prod_{p^a\parallel n}p^{a-2}(p-1)^2
 \ge\frac12:
\]
each factor for $p\ge3$ is at least one, and the sole possible
factor less than one is $1/2$ for $2^1$. Consequently $n\le2D^2$,
so $n\mid N$. This proves (C44.A03.2), including its equivalence
to absence of roots of this form. For $D=0$ the conclusion is
immediate.\hfill$\square$

\paragraph{Why H9-T5 does not supply this test.}
JAC1964 H9-T5 is a theorem about the \emph{odd} variable local
system $R^{2m-1}$ of the pencil and its variable incidence abelian
factor. Its irreducibility hypothesis turns a nonzero injection
$H^1(D)(-(m-1))$ into an isomorphism; one slope outside
$[m-1,m]$ then contradicts that isomorphism. The even-degree
fixed-axis space in (C44.A03.1) is a different representation.
At a closed point the arithmetic Galois group is procyclic, and
an isolated algebraic class need not extend to a flat algebraic
class on the generic axis family. No family-monodromy
irreducibility statement can therefore be substituted into
(C44.A03.1) without a separate spreading theorem.
An out-of-range slope in some other direct summand is compatible
with a slope-$r$ Tate summand. The correct stronger sufficient
Newton condition here is absence of \emph{every} slope-$r$
subspace, while the cyclotomic criterion above is sharper.
This repairs the proposed use of H9-T5 without producing a
witness for every special axis.

\paragraph{Supersingular ambient varieties do not resolve this distinction.}
If $X$ is supersingular, $H^{2r}(X)$ has only slope $r$ because
$H^1(X)$ has slope $1/2$ and the cohomology algebra is exterior.
That identifies the slopes of the ambient summand in
(C44.A01.3), not the slopes or cycle classes of $V_B$.
A nonliftable algebraic class pulled back from $X$ lies in the
ambient summand and is removed by the projector $1-e$. Thus
ambient nonliftability alone gives no nonliftable element of
$A_B$.

\subsection*{C44-A05. The embedded axis itself always has a characteristic-zero lift}

\textbf{Proposition.}
Let $k=\overline{\mathbf F}_p$, let $X/k$ be an abelian variety,
let $L$ be ample, and let $s_0,s_1\in H^0(X,L)$ have a smooth
codimension-two common zero scheme $B$. There is a mixed-characteristic
discrete valuation ring $R$, with finite residue field contained in $k$,
an abelian scheme $\mathcal X/R$, a relatively ample line bundle
$\mathcal L$, and two sections whose common zero scheme
$\mathcal B/R$ is smooth and projective, such that the geometric
special fibre is the given $(X,L,s_0,s_1,B)$. In particular this
applies to every axis in QH9-axis, including supersingular ambient
abelian varieties. The proposition asserts no lifting of a specified
middle-dimensional algebraic class on $B$.

\textbf{Proof.}
Descend the given data to a finite subfield of $k$. The polarized
lifting theorem of Mumford--Norman--Oort gives an abelian scheme
with polarization lifting $(X,\lambda_L)$ over an integral
mixed-characteristic base. An exact statement is Oort,
\emph{Did Earlier Thoughts Inspire Grothendieck?}, Theorem (8.3).1,
p.~29. By finite-presentation descent we may use an integral base
$S$ of finite type over $\mathbf Z$, with a closed point of finite
residue field giving the prescribed geometric special fibre. The
scheme, group laws, and polarization all descend; passing to a
sufficiently large finite-type model preserves smoothness and
properness. This is the usual limit construction, as in Stacks
Project, Lemma 32.10.1, tag 01ZM, Lemma 32.8.9, tag 0C0C,
and Lemma 32.13.1, tag 081F.

The chosen line bundle also lifts after an \'etale neighborhood of
the special point: Olsson, \emph{Compactifications of moduli of
abelian varieties: an introduction}, Lemma 2.1.8 and its proof,
pp.~299--300. Its proof preserves the chosen special-fibre line
bundle. Alternatively, after choosing a representative of the
lifted polarization, the difference from $L$ lies in the dual
abelian variety and lifts after a smooth base change. Take an
integral component through the chosen point and still dominating
$S$; this is possible by flatness of the \'etale neighborhood.
Write $\pi:\mathcal X\to S$ for the resulting family and
$\mathcal L$ for its ample line bundle. For an ample line bundle
on an abelian scheme,
\[
 R^i\pi_*\mathcal L=0\quad(i>0),\qquad
 E=\pi_*\mathcal L\ \text{is locally free and commutes with base change}.
\]
The exact family statement is Olsson, \S2.1.4, p.~299.

The total space $W$ of the vector bundle $E\oplus E$ carries a
universal pair of sections on $\mathcal X\times_S W$. It is
integral of characteristic zero and of finite type over
$\mathbf Z$. The given pair determines a closed point $b\in W$
after a finite residue-field extension. Let $\mathcal Z$ denote
the universal common zero scheme. On $\mathcal Z$, failure of the
two relative differentials to have rank two is closed. Its image
in $W$ is closed because $\mathcal Z\to W$ is proper. The
complement $U$ contains $b$, since $B$ is smooth of codimension
two. The relative Jacobian criterion makes $\mathcal Z_U\to U$
smooth of relative dimension $\dim X-2$, and it is projective.
Thus $U$ is an integral finite-type mixed-characteristic
parameter scheme carrying a smooth projective lift of the
specified geometric fibre.

For completeness one can pass to an arithmetic trait without
changing that geometric fibre. Let $A=\mathcal O_{U,b}$, with
maximal ideal $\mathfrak m$, and put $d=\dim A$. This is an
excellent local domain, $0\ne p\in\mathfrak m$, and its residue
field is finite. Choose $x_1,\ldots,x_{d-1}$ so that
$(p,x_1,\ldots,x_{d-1})$ is $\mathfrak m$-primary. Set
$I=(x_1,\ldots,x_{d-1})$. Some minimal prime $P$ of $I$ does
not contain $p$: otherwise $p\in\sqrt I$ would give
$\sqrt I=\mathfrak m$, contrary to Krull's height theorem.
The local domain $A/P$ has dimension one, since $p$ is nonzero
and generates an ideal with maximal radical in that quotient.
Its normalization is finite by excellence; localization at a
point above its closed point is a discrete valuation ring $R$
with characteristic-zero fraction field and finite residue field.
Pulling back the family along
$\operatorname{Spec}R\to\operatorname{Spec}A\to U$ proves the
claim. The residue extension embeds in $k$, so the geometric
special fibre is the prescribed one.\hfill$\square$

\paragraph{Programme location and remaining obstruction.}
This completes the geometric lifting substep of QH9-axis(a),
in the schemes-over-$\mathbf Z$ framework. The unresolved issue
is surjectivity, for some such lift, of specialization on the
specified rational algebraic middle classes. Smooth proper base
change is an isomorphism on cohomology; it is not a lifting
theorem for cycles. The sign transfer to classes which do lift
is exactly the already settled H9-T8a specialization argument.
Cote 44, Th.~5, scan 004, and the primitive margin on scan 008
therefore retain the same cycle-level gate.

\section{Unresolved axis steps after the ordered attacks}

\paragraph{Lifting.}
For any specified smooth projective mixed-characteristic lift
of $(B,h_B)$ and any rational algebraic class on the geometric
generic fibre specializing to $w\in A_B$, compatibility of
cycle classes with smooth proper specialization carries
$L_Bw=0$ back to the same primitive relation on the generic
fibre. The specialization is an isomorphism on $\ell$-adic
cohomology, so a nonzero $w$ has nonzero generic class, and the
intersection number is unchanged. H9-T8a then gives
$(-1)^r\int_Bw^2>0$. This is the carried lifting argument,
not a surjectivity theorem for Chow specialization.

No proof that every $w\in A_B$ lifts with $B$ is obtained.
Nor is a specific cohomological class in $A_B$ proved
nonliftable in this contribution. Attempts using nonliftable
ambient Tate classes fail because their projection to $A_B$
is zero. Agugliaro, \emph{Examples for the standard conjecture
of Hodge type}, Doc. Math. 31 (2026), Proposition 5.12,
pp.~876--877, supplies a rank-two nonliftable Tate block for
CM lifts of certain abelian varieties; Remark 3.5, p.~870,
explains the passage to arbitrary abelian lifts. Even after
032 supplies algebraicity, these remain ambient examples.
They are not examples for the axis kernel.

\paragraph{Correspondences.}
C44-A02 and JAC1964 H9-T13 eliminate graph/ambient transfers.
A general correspondence from $B$ to an abelian variety could
still act nontrivially, but a splitting preserving the signed
geometric pairing has not been constructed. Merely splitting
an abstract Tate class from $h(B)$ would not identify that
pairing with a positive form.

\paragraph{Newton slopes.}
C44-A03 gives the correct finite-field vanishing test.
No Frobenius polynomial satisfying that test has been exhibited
for every special axis, and axes with $A_B\ne0$ necessarily
fail it. A slope witness in the odd pencil variation does not
settle the even fixed-axis sign.

\paragraph{Current gate.}
The unresolved assertion remains weak nonnegativity of
$(-1)^r\int_Bw^2$ for every algebraic $w\in\ker\kappa_*$.
Conditional on 032/Milne, C44-A01 identifies it with the weak
Hodge-standard assertion for $B$. No wrong-sign cycle is
constructed. No main question is marked closed by these
route boundaries.
